\chapter{Decomposition, abstraction, and boundaries}

\section{The purpose of decomposition}
Decomposition turns a large problem into units that people can understand,
change, and test. It is not merely splitting a file. A good boundary hides a
decision that is likely to change or that should be protected by an invariant.
Parnas's classic discussion treats information hiding as a criterion for module
design, rather than a preference for a particular syntax \cite{parnas1972criteria}.

For Ledgerline, a useful first decomposition is:

\begin{center}
\begin{tikzpicture}[node distance=8mm and 12mm,>=Latex]
  \node[draw,rounded corners,fill=codebg,minimum width=30mm,minimum height=9mm] (api) {HTTP adapter};
  \node[draw,rounded corners,fill=codebg,minimum width=30mm,minimum height=9mm,below=of api] (app) {application service};
  \node[draw,rounded corners,fill=codebg,minimum width=30mm,minimum height=9mm,below=of app] (domain) {order domain};
  \node[draw,rounded corners,fill=codebg,minimum width=30mm,minimum height=9mm,left=of domain] (store) {store port};
  \node[draw,rounded corners,fill=codebg,minimum width=30mm,minimum height=9mm,right=of domain] (events) {event port};
  \node[draw,rounded corners,fill=codebg,minimum width=30mm,minimum height=9mm,below=of store] (sql) {SQL adapter};
  \node[draw,rounded corners,fill=codebg,minimum width=30mm,minimum height=9mm,below=of events] (queue) {queue adapter};
  \draw[->,thick,accent] (api) -- (app);
  \draw[->,thick,accent] (app) -- (domain);
  \draw[->,thick,accent] (domain) -- (store);
  \draw[->,thick,accent] (domain) -- (events);
  \draw[->,thick,accent] (store) -- (sql);
  \draw[->,thick,accent] (events) -- (queue);
\end{tikzpicture}
\end{center}

The arrows represent dependency direction, not necessarily runtime network
traffic. The domain does not need to know whether persistence is SQLite,
PostgreSQL, or an in-memory fake. The adapter owns translation between a
technology's vocabulary and the domain's vocabulary.

\section{Cohesion and coupling}
\term{Cohesion} asks whether the responsibilities inside a unit belong together.
\term{Coupling} asks how much one unit knows about another. High cohesion and
low unnecessary coupling are useful heuristics, not scores to maximize blindly.
Some coupling is valuable: an invariant about money may be kept close to the
database constraint that enforces it. The danger is accidental coupling, such as
a UI importing a database row shape and thereby making a storage migration a UI
change.

A dependency is especially expensive when it crosses:

\begin{itemize}
  \item ownership boundaries, because coordination becomes necessary;
  \item time boundaries, because queues and retries introduce duplication or delay;
  \item trust boundaries, because input must be authenticated and validated;
  \item data-lifetime boundaries, because privacy and deletion rules differ;
  \item release boundaries, because two versions must coexist.
\end{itemize}

\section{Abstraction and the cost of being wrong}
An abstraction removes detail so that a person can reason about a larger unit.
It also hides information. An abstraction is successful when the hidden detail
can vary without surprising its clients. It is unsuccessful when the abstraction
promises a simple concept but leaks the complexity through exceptions, latency,
ordering, or resource usage.

The phrase “repository” can hide a database, but it cannot make a database
transaction disappear. If callers assume that `save` is atomic while the
implementation performs three independent writes, the abstraction lies. Make
important semantics explicit: \texttt{save\_order\_and\_outbox\_event} communicates more
than `save` if atomicity is essential.

\section{Designing a boundary}
For each proposed module, write four sentences:

\begin{enumerate}
  \item Its responsibility is ...
  \item It guarantees ...
  \item It may change without requiring ...
  \item It must not know ...
\end{enumerate}

Then test the boundary against a plausible change. If changing the tax rate
requires editing HTTP parsing, SQL queries, and every unit test, the domain rule
is scattered. If a tax module cannot be tested without a network and a database,
the boundary is probably too entangled.

\section{When a boundary is too early}
Premature decomposition creates indirection before the shape of the problem is
understood. Five interfaces around five lines of code can make a small system
harder to read. Start with seams that protect a real decision: an external API,
a security check, a durable write, or a volatile policy. Refactor toward clearer
boundaries as evidence accumulates. The repository's Python example keeps the
domain function small because the purpose is visible; it does not manufacture a
framework around it.

\begin{exercise}
Take a single function that validates and persists an order. List three changes
that could occur independently. Design boundaries that isolate those changes,
then name one cost introduced by each boundary.
\end{exercise}

\section{Chapter summary}
Decompose around changing decisions and protected invariants. Use information
hiding, cohesion, and coupling as reasoning tools. Treat abstractions as
contracts about semantics, not just shorter names, and delay indirection when it
does not yet protect a real change.
